% !PRINTED 619
\weberpage{619}
%<<CONT
die Ziffern von $c$ ungeändert, während sie doch nicht die identische Substitution ist, weil $c_1^h$ nicht alle Ziffern ungeändert lässt. Dies ist aber nicht möglich, wenn $h>1$ ist. Es muss also entweder $\pi$ eine Ziffer ungeändert lassen, oder es muss aus Cyklen von gleichem Grade bestehen. Da aber $n$ Primzahl ist, so kann $\pi$ in diesem Falle nur einen Cyklus vom $n^{\mathrm{ten}}$ Grade bilden.

Es enthält also $P$ ausser der identischen nur cyklische Permutationen $n^{\mathrm{ten}}$ Grades, die wir mit $\gamma$ bezeichnen wollen, und Permutationen $\varkappa$, die eine Ziffer ungeändert lassen und ausserdem aus Cyklen von gleichem Grade bestehen.

Jede der Permutationen $\varkappa$, durch die die Ziffer 0 ungeändert bleibt, wollen wir mit $\varkappa_0$ bezeichnen. Ebenso bedeuten $\varkappa_1, \varkappa_2 \ldots \varkappa_{n-1}$ die Permutationen $\varkappa$, die die Ziffern 1, 2 \ldots $n-1$ ungeändert lassen. Da $P$ transitiv ist, so müssen ebensoviel $\varkappa_0$, wie $\varkappa_1$, wie $\varkappa_2$ etc. vorhanden sein.

Denn ist $\pi$ eine Permutation, die 0 in 1 überführt, so ist jedes $\pi^{-1}\varkappa_0\pi$ ein $\varkappa_1$, und umgekehrt jedes $\pi\varkappa_1\pi^{-1}$ ein $\varkappa_0$, und ebenso für die übrigen Ziffern. Ist also $\mu$ die Anzahl der $\varkappa_0$, $\nu$ die Anzahl der $\gamma$, $m$ der Grad von $P$, so ist, da noch die identische Permutation hinzukommt,

\weq{11}{m = \mu n + \nu + 1.}

Nun bilden die $\varkappa_0$ mit der identischen Permutation zusammen eine Gruppe $Q$ vom Grade $\mu+1$ und einen Theiler von $P$, nämlich den Inbegriff aller Permutationen von $P$, die 0 ungeändert lassen. Sind also $\pi_1, \pi_2 \ldots \pi_{n-1}$ Permutationen in $P$, die 0 in $1, 2 \ldots n-1$ überführen, so ist $Q\pi_1$ der Inbegriff aller der Permutationen von $P$, die 0 in 1 überführen u. s. f.

Wir erhalten also die Zerlegung von $P$ in die Nebengruppen

\[
P = Q + Q\pi_1 + Q\pi_2 + \cdots + Q\pi_{n-1},
\]

woraus folgt

\weq{12}{m = n(\mu+1),}

also aus (11) und (12) $\nu = n-1$.

Es giebt also $n-1$ und nicht mehr cyklische Permutationen $n^{\mathrm{ten}}$ Grades in $P$ und diese bilden folglich wieder mit der Einheit zusammen eine cyklische Gruppe

\[
C = 1, \gamma, \gamma^2 \ldots \gamma^{n-1},
\]

da mit $\gamma$ zugleich alle Potenzen von $\gamma$ in $P$ vorkommen.

Wenn nun $\gamma$ eine cyklische Permutation $n^{\mathrm{ten}}$ Grades ist, so ist auch jedes $\pi^{-1}\gamma^h\pi$ cyklisch und muss also, wenn $\pi$ zu $P$
%CONT>>
